\chapter{シェルスクリプトによる自動化と\mbox{一括処理}}
\label{app:shell_scripting}
\index{しぇるすくりぷと@シェルスクリプト|textbf}

第~\ref{chap:shell}~章では、シェルを対話的に操作する方法を説明した。同じコマンド列を複数のファイルや条件に対して実行する場合は、操作をシェルスクリプトへ記述すると、処理順序と引数を記録したまま再実行できる。

本付録では、引数の受取り、変数の計算、ファイルの存在確認、条件分岐、配列、\code{for} ループを扱う。最後に、設定ファイルのテンプレートを \code{sed} で置換して複数条件を実行する方法を説明する。

\section{シェルスクリプトの基礎構造}

\subsection{シバンと実行権限}
\index{しばん@シバン|textbf}

シェルスクリプトの先頭には、どのシェルで解釈するかを指定するシバン（shebang）を書く。本付録の例は Bash の配列と算術式を使うため、\code{env} で \code{bash} の所在を検索する。

\begin{lstlisting}[language=bash]
#!/usr/bin/env bash

echo "Hello, Shell Scripting!"
\end{lstlisting}

\code{run.sh} という名前で保存した場合は、実行権限を付けてから直接実行する。

\begin{lstlisting}[numbers=none, xleftmargin=2\zw]
$ chmod +x run.sh
$ ./run.sh
Hello, Shell Scripting!
\end{lstlisting}

\code{chmod +x} は実行権限を追加し、\code{./run.sh} は現在のディレクトリにあるファイルを実行する。最後の行がスクリプトの標準出力である。

\subsection{コマンドライン引数}
\index{こまんどらいんひきすう@コマンドライン引数|textbf}

スクリプトに外から値（ファイル名や設定値）を渡すには、実行時の引数を利用する。スクリプト内からは \code{\$1}, \code{\$2} といった特殊な変数で受け取れる。

\begin{itemize}
    \item \code{\$0}: 実行されたスクリプトのファイル名（例: \code{./run.sh}）
    \item \code{\$1, \$2, ...}: 1 番目、2 番目以降の引数
    \item \code{\$\#}: 渡された引数の総数
    \item \code{\$@}: 渡された全ての引数
\end{itemize}

次のスクリプトは、二つの位置引数を検査してから入力名と出力名を表示する。

\begin{lstlisting}[language=bash]
#!/bin/bash

# 引数が2つ渡されたかどうかを確認する簡易チェック
if [ "$#" -ne 2 ]; then
    echo "Usage: $0 <input_file> <output_file>"
    exit 1
fi

INPUT=$1
OUTPUT=$2
echo "Reading from $INPUT, writing to $OUTPUT"
\end{lstlisting}

\code{./run.sh input.csv result.csv} と実行すると、\code{Reading from input.csv, writing to result.csv} と表示される。引数が二つでなければ使用法を表示し、終了状態 1 で処理を打ち切る。

\section{変数、演算、条件式}

\subsection{変数の定義と配列}
\index{bash@Bash!はいれつ@配列}

シェルで変数を定義する際は、イコールの前後に空白を入れてはならない。また、変数の値を参照するときは先頭に \code{\$} を付ける。

\begin{lstlisting}[language=bash]
COUNT=100
NAME="data_analysis"

echo "The count is $COUNT."       # 展開される
echo 'The count is $COUNT.'       # シングルクォートでは展開されず "$COUNT" と表示される
\end{lstlisting}

Bash では、要素を空白で区切って括弧 \code{()} で囲むと配列を定義できる。各要素を引用符で囲めば、要素内の空白やワイルドカードが意図せず展開されることを防げる。

\begin{lstlisting}[language=bash]
# 配列の定義
targets=("background" "signal" "noise")

# 配列サイズの取得
echo "Number of targets: ${#targets[@]}"

# 配列要素へのアクセス
echo "First element: ${targets[0]}"

# 配列全体を展開する（for ループでよく使う）
echo "All elements: ${targets[@]}"
\end{lstlisting}

\subsection{数値計算（算術式展開）}
\index{さんじゅつしきてんかい@算術式展開|textbf}

Bash で整数の加減乗除を行う場合は、算術式展開 \code{\$((...))} または算術評価 \code{((...))} を用いる。この中では変数名の \code{\$} を省略できる。Bash の算術評価は浮動小数点計算に対応しないため、小数を扱う場合は Python、\code{bc}、\code{awk} などを用いる。

\begin{lstlisting}[language=bash]
A=10
B=5

# 計算式は (( )) を用いる
SUM=$((A + B))
MUL=$((A * B))

# 変数を直接加算・インクリメントする
(( A++ ))
(( B += 20 ))
\end{lstlisting}

\subsection{条件式（\code{test} コマンド）}
\index{test command@\code{test}|textbf}

数値の大小比較や、ファイルが存在するかどうかの確認には、\code{[ ]} を用いる。大括弧の内側には必ずスペースが必要である。

\textbf{数値の比較:}
\begin{itemize}
    \item \code{-eq}: 等しい (equal)
    \item \code{-ne}: 等しくない (not equal)
    \item \code{-gt}: より大きい (greater than)
    \item \code{-ge}: 以上 (greater than or equal)
    \item \code{-lt}: 未満 (less than)
    \item \code{-le}: 以下 (less than or equal)
\end{itemize}

\textbf{文字列の比較:}
\begin{itemize}
    \item \code{=}, \code{==}: 一致する
    \item \code{!=}: 一致しない
    \item \code{-z "\$STR"}: 文字列が空である
\end{itemize}

\textbf{ファイルやディレクトリの確認:}
\begin{itemize}
    \item \code{-e "\$FILE"}: ファイル（またはディレクトリ）が存在する
    \item \code{-f "\$FILE"}: 通常のファイルとして存在する
    \item \code{-d "\$DIR"}: ディレクトリとして存在する
\end{itemize}

これらの条件式は、後述する \code{if} 文の \code{[]} の中身として頻繁に登場する。

\section{条件分岐 (\code{if}, \code{case})}

\subsection{\code{if} 文を用いた分岐}
\index{bash@Bash!if@\code{if} 文}

前節の条件式を組み合わせて処理を分岐する。\code{then} と \code{fi} でブロックを囲むのがシェルの特徴である。

\begin{lstlisting}[language=bash]
#!/bin/bash
FILE=$1

if [ -z "$FILE" ]; then
    echo "エラー: ファイル名が指定されていません"
    exit 1
fi

if [ -f "$FILE" ]; then
    echo "ファイル $FILE を発見しました。処理を開始します。"
elif [ -d "$FILE" ]; then
    echo "エラー: $FILE はディレクトリです。ファイルを指定してください。"
    exit 1
else
    echo "エラー: $FILE が存在しません。"
    exit 1
fi
\end{lstlisting}

\subsection{\code{case} 文による分岐}
\index{bash@Bash!case@\code{case} 文}

条件が特定の文字列パターンとの一致で表せる場合は、\code{case} 文を用いると \code{elif} を重ねずに分岐を列挙できる。拡張子の判定や、\code{-h}、\code{--version} などのオプション処理に利用できる。

\begin{lstlisting}[language=bash]
#!/bin/bash
MODE=$1

case "$MODE" in
    "start")
        echo "Starting process..."
        ;;
    "stop" | "halt")
        echo "Stopping process..."
        ;;
    *)
        echo "Unknown mode: $MODE"
        echo "Available modes: start, stop, halt"
        exit 1
        ;;
esac
\end{lstlisting}

各区分の終わりは \verb|;;| で締めくくる。アスタリスク \code{*} は「上記以外のすべて」にマッチするワイルドカードとして機能する。

\section{反復処理と一括自動化 (\code{for})}

\subsection{配列やリストの要素に対する \code{for} ループ}
\index{bash@Bash!for@\code{for} ループ}

指定した値のリストや、配列の全要素に対して処理を繰り返す \code{for} 文の例を次に示す。

\begin{lstlisting}[language=bash]
#!/bin/bash

# 配列の全要素に対して
targets=("run_01" "run_02" "run_03")

for t in "${targets[@]}"; do
    echo "Processing target: $t"
done

# 単純な数値の手書きリストやコマンド展開に対して
for num in 10 20 30; do
    echo "Number is $num"
done
\end{lstlisting}

\subsection{実践: 複数ファイルのパス置換とバッチ処理}
\index{ばっちしょり@バッチ処理|textbf}

実務で最もよく使うループは「特定のディレクトリにあるファイルを順次処理する」ための構文である。ここで、入力ファイル名（例: \code{data/input/run01.csv}）から、出力ファイル名（例: \code{data/output/run01\_result.png}）を自動的に作り出す技術が必要となる。

\code{basename} コマンドや、変数の後方一致削除（\code{\%...}）を用いることでこれを簡潔に実装できる。

\begin{lstlisting}[language=bash]
#!/bin/bash

INPUT_DIR="data/input"
OUTPUT_DIR="data/output"

# 出力先が存在しなければ作成する
if [ ! -d "$OUTPUT_DIR" ]; then
    mkdir -p "$OUTPUT_DIR"
fi

# INPUT_DIR の中にあるすべての .csv ファイルに対して反復
for filepath in "$INPUT_DIR"/*.csv; do
    # 該当するファイルが存在しない場合はそのままループが回ってしまうのでチェックする
    if [ ! -f "$filepath" ]; then
        echo "No CSV files found in $INPUT_DIR"
        break
    fi

    echo "Found raw file: $filepath"

    # basename コマンドでディレクトリパスを取り除き、ファイル名だけにする
    # 例: "data/input/run01.csv" -> "run01.csv"
    filename=$(basename "$filepath")

    # 文字列置換: 後方の ".csv" を取り除く
    # 例: "run01.csv" -> "run01"
    name_no_ext="${filename%.csv}"

    # 新しい出力パスを組み立て直す
    output_path="$OUTPUT_DIR/${name_no_ext}_result.png"

    echo "Running analysis Python code for: $name_no_ext"
    # 実際の処理コマンド（例: python analyze.py <in> <out>）
    # python analyze.py "$filepath" "$output_path"
done
\end{lstlisting}

この例では、入力ファイルごとに拡張子を除いた名前を作り、対応する出力パスを組み立てている。解析コマンドを有効にする前に、\code{echo} で表示される入出力の対応を確認すれば、意図しない上書きを防げる。

\section{\code{sed} を用いたテンプレート書き換えと動的実行}
\index{sed@\code{sed}|textbf}

外部の実行プログラムの中には、引数（\code{--input XXX}）の形式ではなく、設定ファイル（例: \code{config.ini}）を読み込ませることで動作するものがある。バッチ処理で繰り返し実行したい場合、ループを回すたびに設定ファイルの特定のパラメータ（例えばシード値や閾値など）を書き換える必要が生じる。

\code{sed} を使うと、テンプレートとして用意した設定ファイルの一部を置換し、条件ごとの実行用設定ファイルを生成できる。

\subsection{テンプレートファイルの準備}

テンプレートでは、書き換えるパラメータをファイル内の他の場所に現れないプレースホルダーへ置き換える。次の例では、\verb|__THRESHOLD__| と \verb|__OUTFILENAME__| を用いる。

\begin{lstlisting}[numbers=none, title={config\_template.ini}]
[Simulation]
Events = 10000
Threshold = __THRESHOLD__
OutputFile = __OUTFILENAME__
\end{lstlisting}

\subsection{\code{sed} による置換の実行}

シェルスクリプト内で条件値を順に取り出し、\code{sed} で置換した設定ファイルを生成して実行プログラムへ渡す。

基本形は \code{sed 's/置換前文字列/置換後文字列/g' 入力ファイル > 出力ファイル} である。\code{sed} の置換式でシェル変数を展開するには、シングルクォートではなく\textbf{ダブルクォート（\code{"..."}）}で囲む必要がある。

\begin{lstlisting}[language=bash]
#!/bin/bash

# 閾値を 3 つ変えてシミュレーション用の設定ファイルを動的生成し、実行する例
thresholds=(0.5 1.0 2.5)

for thresh in "${thresholds[@]}"; do

    out_name="result_th_${thresh}.dat"
    run_config="config_run.ini"

    echo "Generating config for Threshold = $thresh"

    # sedコマンドを用いて __THRESHOLD__ と __OUTFILENAME__ の 2 箇所を同時に置換する。
    # 変数 ($thresh など) を展開させるためダブルクォートで 's/.../.../g' 全体を囲む
    sed -e "s/__THRESHOLD__/$thresh/g" \
        -e "s/__OUTFILENAME__/$out_name/g" \
        config_template.ini > "$run_config"

    # 置換された設定の内容を出力して確認 (または直接プログラムを走らせる)
    cat "$run_config"

    # 生成した設定ファイルを実行プログラムへ渡す
    # ./my_simulation_app --config "$run_config"

done
\end{lstlisting}

この方法は、設定ファイルを入力とする Python プログラムだけでなく、コンパイル済みのシミュレーションプログラムにも適用できる。各反復で生成した設定ファイルを別名で保存すれば、どの条件で各結果を得たかを後から確認できる。ただし、置換値に \code{/} や \code{\&} が含まれる場合は、\code{sed} の区切り文字やエスケープ方法を調整する必要がある。
